% ===== PRÉAMBULE CANONIQUE — à coller VERBATIM en tête de CHAQUE .tex =====
\documentclass[11pt,a4paper]{article}
\usepackage[utf8]{inputenc}\usepackage[T1]{fontenc}\usepackage{lmodern}
\usepackage[french]{babel}
\usepackage{amsmath,amssymb,mathtools}\usepackage{icomma}
\usepackage{siunitx}\sisetup{locale=FR}  % \qty{}{}, \si{}, \ang{}
\usepackage{xcolor}\usepackage{colortbl}
\usepackage{tikz}\usepackage{pgfplots}\pgfplotsset{compat=1.18}
\usepackage[siunitx,europeanresistors]{circuitikz} % circuits (élec.)
\usepackage[most]{tcolorbox}\usepackage{enumitem}\usepackage{array,booktabs}
\usepackage[a4paper,margin=2.2cm]{geometry}
\usetikzlibrary{arrows.meta,calc,angles,quotes,positioning,patterns,%
  backgrounds,shapes.geometric,decorations.pathmorphing,decorations.markings,intersections}
\definecolor{primary}{RGB}{222,49,99}\definecolor{primarydark}{RGB}{150,22,55}
\definecolor{defcol}{RGB}{0,114,178}\definecolor{methcol}{RGB}{52,92,148}
\definecolor{excol}{RGB}{0,135,90}\definecolor{attncol}{RGB}{200,40,55}
\tcbset{boxbase/.style={enhanced,breakable,boxrule=0.9pt,arc=2.5pt,
  left=3mm,right=3mm,top=2.3mm,bottom=2.3mm,fonttitle=\bfseries,coltitle=white}}
% --- boîtes TUTORIEL (noms EXACTS, ne pas renommer) ---
\newtcolorbox{cle}[1][]{boxbase,colback=defcol!6,colframe=defcol,
  title={\textsc{À retenir}\ifx\relax#1\relax\relax\else\textmd{ : \itshape #1}\fi}}
\newtcolorbox{methode}[1][]{boxbase,colback=methcol!7,colframe=methcol,sharp corners,
  title={\textsc{Méthode}\ifx\relax#1\relax\relax\else\textmd{ --- \itshape #1}\fi}}
\newtcolorbox{exemplebox}[1][]{boxbase,colback=excol!5,colframe=excol,
  title={\textsc{Exemple}\ifx\relax#1\relax\relax\else\textmd{ : \itshape #1}\fi}}
\newtcolorbox{piege}[1][]{boxbase,colback=attncol!6,colframe=attncol,
  title={\textsc{Piège}\ifx\relax#1\relax\relax\else\textmd{ : \itshape #1}\fi}}
\newtcolorbox{remarque}[1][]{boxbase,colback=black!4,colframe=black!45,
  title={\textsc{Remarque}\ifx\relax#1\relax\relax\else\textmd{ : \itshape #1}\fi}}
% --- boîtes EXERCICE / CORRIGÉ (noms EXACTS, auto-comptées) ---
\newtcolorbox[auto counter]{exo}[1][]{boxbase,colback=primary!4,colframe=primary,
  title={\textsc{Exercice}~\thetcbcounter\ifx\relax#1\relax\relax\else\textmd{ --- \itshape #1}\fi}}
\newtcolorbox[auto counter]{corrige}[1][]{boxbase,colback=excol!4,colframe=excol,
  title={\textsc{Corrigé}~\thetcbcounter\ifx\relax#1\relax\relax\else\textmd{ --- \itshape #1}\fi}}
\newcommand{\vv}[1]{\vec{#1}}\newcommand{\ds}{\displaystyle}
\newcommand{\R}{\mathbb{R}}\newcommand{\N}{\mathbb{N}}\newcommand{\Z}{\mathbb{Z}}
% =========================================================================
\begin{document}
\sisetup{per-mode=symbol}

\begin{exo}[la pierre au bout d'un fil]
Une pierre de masse $m=\qty{0,2}{\kilogram}$ tourne au bout d'un fil de longueur $R=\qty{0,5}{\metre}$ et
fait un tour en $T=\qty{1}{\second}$ ($\pi\approx\num{3,14}$).
\begin{enumerate}[label=\alph*)]
  \item Calcule la vitesse linéaire $v$.
  \item Calcule l'accélération centripète $a_c$.
  \item Calcule la force centripète $F_c$ (fournie par la tension du fil).
\end{enumerate}
\end{exo}

\begin{exo}[la force centripète]
Un objet de masse $m=\qty{2}{\kilogram}$ décrit un cercle de rayon $R=\qty{3}{\metre}$ à la vitesse
linéaire $v=\qty{6}{\metre\per\second}$. Calcule la force centripète $F_c=\dfrac{mv^2}{R}$.
\end{exo}

\begin{exo}[deux enfants sur un manège]
Un manège tourne à la vitesse angulaire $\omega=\qty{2}{\radian\per\second}$. Un enfant est assis à
$R_1=\qty{1}{\metre}$ du centre, l'autre à $R_2=\qty{3}{\metre}$.
\begin{enumerate}[label=\alph*)]
  \item Calcule la vitesse linéaire de chaque enfant.
  \item Lequel se déplace le plus vite ? Ont-ils la même vitesse angulaire ?
\end{enumerate}
\end{exo}

\begin{exo}[retrouver la période]
Un satellite parcourt un cercle de rayon $R=\qty{2}{\metre}$ (modèle réduit) à la vitesse linéaire
$v=\qty{4}{\metre\per\second}$ ($\pi\approx\num{3,14}$).
\begin{enumerate}[label=\alph*)]
  \item À partir de $v=\dfrac{2\pi R}{T}$, calcule la période $T$.
  \item En déduis la fréquence $f$.
\end{enumerate}
\end{exo}

\begin{exo}[doubler la vitesse]
Un objet en MCU subit une force centripète $F_c=\qty{10}{\newton}$. On garde le même rayon $R$ et la
même masse, mais on \emph{double} sa vitesse linéaire.
\begin{enumerate}[label=\alph*)]
  \item Par quel facteur la force centripète est-elle multipliée ?
  \item Que vaut la nouvelle force centripète ?
\end{enumerate}
\end{exo}
\end{document}
